\subsection{A binary tetrahedral lifting criterion}
\label{sec:163-binary-tetrahedral}

The remaining central obstruction has an exact finite lifting
interpretation. The criterion below is conditional on the larger
scalar image and does not establish the existence of the lift.

Write
\[
 G=\widehat{\mathrm{GT}},\quad H=G_{\mathbf Q},\quad
 R=R_{\rm global},\quad \Gamma=G/R,
 \qquad G=RH.
\]
Suppose that the larger degree-163 image in
Proposition~\ref{prop:cm163-full-image-dichotomy} occurs.  Thus
\[
 h:R\twoheadrightarrow P,\qquad h(R\cap H)=A,\qquad
 E=P/A\simeq\mathbf F_2^2,
 \qquad J=\ker h.
\]
The cyclotomic action on \(E\) factors through a surjection
\(\gamma:\Gamma\to C_3\) and is faithful and irreducible.
Proposition~\ref{prop:163-relative-cohomology} identifies the invariant
relative degree-two group with the line generated by the pullback of
\(Q_8\to E\). Theorem~\ref{thm:163-concrete-odd-inertia} gives
\[
 H^1_{\rm cont}(\Gamma,E^\vee)=0.
\]
The polarization of \(q(z)=z^3\), namely
\(\langle z,w\rangle=zw^2+z^2w\), identifies \(E\) with
\(E^\vee\) equivariantly.  Therefore
\begin{equation}
 H^1_{\rm cont}(\Gamma,E)=0.                 \label{eq:163-selfdual-h1-vanishing}
\end{equation}

The normalized relative degree--163 cocycle and its finite quotient are
\[
 z(rh_0)=h(r)\bmod A,\qquad
 \pi(g)=(z(g),\gamma(g)):
 G\twoheadrightarrow E\rtimes C_3\simeq A_4.
\]
Here \(h_0\in H\); the symbol \(h\) without a subscript denotes
the scalar homomorphism.  In particular \(z|_H=0\).
Let
\[
 \widetilde A_4=Q_8\rtimes C_3,
 \qquad p:\widetilde A_4\to E\rtimes C_3
\]
where the chosen order-three automorphism cycles the three quaternion
directions.  This is the binary tetrahedral group
\(\mathrm{SL}_2(\mathbf F_3)\). Explicitly, in the latter group
\[
 i=\begin{pmatrix}0&1\\-1&0\end{pmatrix},\qquad
 j=\begin{pmatrix}1&1\\1&-1\end{pmatrix}
\]
satisfy \(i^2=j^2=-I\) and \(ij=-ji\), so they generate
\(Q_8\). An element of order three acts cyclically on its three
noncentral pairs. This gives the semidirect product of order
\(8\cdot3=24=|\mathrm{SL}_2(\mathbf F_3)|\).  Its center is the kernel
\(C_2\), and its displayed \(C_3\) maps isomorphically to the
displayed complement in \(A_4\).

\begin{proposition}\label{prop:163-binary-tetrahedral-lift}
The following conditions are equivalent.
\begin{enumerate}
\item Some continuous, \(G\)-invariant
\(\kappa:J\to\mathbf F_2\), zero on \(J\cap H\), has
nonzero central relative obstruction \(o_2(\kappa)\).
\item There is a continuous homomorphism
\(\widetilde\pi:G\to\widetilde A_4\) such that
\(p\widetilde\pi=\pi\) and
\(\widetilde\pi(H)\) lies in the displayed \(C_3\).
\item The pullback of
\(\widetilde A_4\to A_4\) by \(\pi\), equipped with its
canonical splitting over \(H\), is split as a relative central
extension of \((G,H)\).
\end{enumerate}
Every homomorphism in condition two is surjective and has
\(\widetilde\pi(H)=C_3\).  When such homomorphisms exist, their
set is a torsor under the continuous quadratic characters of \(G\)
which vanish on \(H\).
\end{proposition}

\begin{proof}
Assume condition one and put \(K=\ker\kappa\).  Invariance makes
\(K\) a closed normal subgroup of \(G\).  Since
\(o_2(\kappa)\ne0\), Proposition~\ref{prop:163-relative-cohomology}
gives an equivalence of central extensions, preserving the specified
arithmetic splitting,
\[
 R/K\simeq P\times_E Q_8.
\]
The image of \(R\cap H\) under this equivalence is
\[
 L=\{(a,1):a\in A\}.
\]
It is central in \(R/K\).  It is also preserved by the image of
\(H\), since it is exactly the image of \(R\cap H\).
As \(G=RH\), it follows that \(L\) is normal in \(G/K\).
All groups and images involved are compact, so \(L\) is closed.

Quotienting \(G/K\) by \(L\) now gives a profinite group
\(G_1\), with subgroups \(R_1\) and \(H_1\) which are the
images of \(R\) and \(H\).  They satisfy
\[
 R_1\simeq Q_8,\qquad R_1\cap H_1=1,\qquad G_1=R_1H_1.
\]
For the intersection, the image of \(H\) in \(G/K\) meets
\(R/K\) in the image of \(R\cap H\), because \(K\subset R\).
That intersection is precisely \(L\), which has just been killed.
Moreover \(H_1\simeq H/(H\cap R)\simeq\Gamma\).  Thus
\[
 G_1\simeq Q_8\rtimes_\alpha\Gamma
\]
for a continuous action \(\alpha:\Gamma\to\operatorname{Aut}(Q_8)\).

The induced action on \(Q_8/Z(Q_8)=E\) is the specified
\(C_3\)-action.  There is an exact sequence
\[
 1\longrightarrow\operatorname{Inn}(Q_8)\simeq E
 \longrightarrow\operatorname{Aut}(Q_8)
 \longrightarrow\operatorname{GL}_2(\mathbf F_2)
 \longrightarrow1.
\]
Indeed an automorphism can send the first two quaternion generators
to arbitrary signed independent directions; there are \(6\cdot4=24\)
choices.  Those inducing the identity on the three directions are
the four inner automorphisms.  Let
\(\theta:C_3\to\operatorname{Aut}(Q_8)\) be the displayed
order-three lift.  The difference between \(\alpha\) and
\(\theta\gamma\) is a continuous cocycle
\(v:\Gamma\to E\), through
\[
 \alpha(\delta)=\operatorname{Inn}(v(\delta))
                         \theta(\gamma(\delta)).
\]
Equation~\eqref{eq:163-selfdual-h1-vanishing} says that \(v\) is a coboundary.  Changing the
identification of \(Q_8\) by one inner automorphism therefore makes
\(\alpha=\theta\gamma\).  This change acts trivially on \(E\).
The map
\[
 Q_8\rtimes_\alpha\Gamma\longrightarrow Q_8\rtimes C_3,
 \qquad(q,\delta)\longmapsto(q,\gamma(\delta))
\]
then gives the required lift of the original \(\pi\), with
the required arithmetic complement.  It is onto.

Conversely, suppose condition two holds.  The image of \(R\)
lies in \(Q_8\), and maps onto \(E\), because \(z(R)=E\).
No proper subgroup of \(Q_8\) maps onto \(E\): every proper
subgroup is cyclic or central.  Hence
\(\widetilde\pi(R)=Q_8\).  Its commutator subgroup is the
central \(C_2\), and \([R,R]\subset J\).  Therefore the
restriction of \(\widetilde\pi\) to \(J\) is a nonzero
character
\[
 \kappa:J\longrightarrow Z(Q_8)=\mathbf F_2.
\]
It is \(G\)-invariant, since this center is central in
\(\widetilde A_4\), and it vanishes on \(J\cap H\), since
the displayed \(C_3\) meets the center trivially.

The homomorphism
\[
 R\longrightarrow P\times_E Q_8,
 \qquad r\longmapsto(h(r),\widetilde\pi(r))
\]
is surjective.  It surjects to \(P\), and its kernel above
\(0\in P\) maps onto the central \(C_2\), by the preceding
nonvanishing of \(\kappa\).  Hence every two-element fiber is
covered.  Its kernel is \(\ker\kappa\), and it sends the
arithmetic intersection to \(\{(a,1):a\in A\}\).
Consequently its relative class is exactly the quaternion class.
This proves condition one.

The equivalence of conditions two and three is the usual elementary
pullback description: a relative splitting is a homomorphic section
\(g\mapsto(g,\widetilde\pi(g))\), whose prescribed value on
\(H\) is the canonical lift through its \(C_3\)-image.
If two such lifts exist, their pointwise ratio lies in the common
central kernel and is a continuous character \(G\to\mathbf F_2\)
vanishing on \(H\).  Conversely multiplying a lift by any such
character gives another lift.  Finally the image of \(H\) is all
\(C_3\), since \(\gamma(H)=C_3\).
\end{proof}

The proposition identifies the finite quotient which would witness
a nonzero central obstruction.  No exclusion or realization of this
binary tetrahedral quotient for full GT has been proved here.  In
particular the result is not a proof that \(o_2\) is zero or nonzero,
does not decide \(h(R)=A\) versus \(P\), and does not decide
\(\operatorname{Hom}_{\rm cont}(B,\mathbf F_2)\).

\paragraph{Commutator evaluation.}
There is also a direct evaluation test.  If \(r_1,r_2\in R\) map to
an \(\mathbf F_2\)-basis of \(E\), then
\[
 o_2(\kappa)\ne0
 \quad\Longleftrightarrow\quad
 \kappa([r_1,r_2])=1.
\]
Indeed the unique invariant nonzero quadratic form has nonzero
polarization on every independent pair.  This test still needs a
character \(\kappa\) or a full-GT identity; it does not
produce either one.
